
\chapter{Эмитент и~эквайер}\label{ch:issuer-acquirer}

Касание карты у~терминала или нажатие \enquote{Оплатить} в~браузере запускает
два встречных процесса. Эмитент за~секунды решает, готов ли взять на~себя
кредитный риск: действительна ли карта, хватает ли средств, не~похожа ли
операция на~мошенничество. Эквайер принимает запрос, проводит его по~нужному
маршруту и~рассчитывается с~продавцом.

\section{Эмитент}\label{sec:ia-issuer}

\highlight{Эмитент (issuer) выпускает карту и~принимает на~себя кредитный риск
транзакции}. В~момент покупки он должен решить, готов ли рассчитаться
по~операции через карточную сеть, а~затем взыскать средства с~держателя
карты~--- списать с~дебетового счёта или увеличить задолженность по~кредитному.
Поэтому у~эмитента работа распределена между специализированными системами:
одна ведёт жизненный цикл карты, другая принимает авторизационное решение,
третья отвечает за~аутентификацию
в~интернет-платежах.

В~клиринг и~расчёт эмитент входит только как получатель: пакет финансовых
требований формирует и~отправляет в~сеть эквайер, а~эмитент по~этому файлу
списывает с~держателя. У~банка, который занимается только эмиссией, клиринговой
инициативы нет~--- он ничего в~сеть не~отправляет, кроме ответов на~чужие
запросы (авторизационные, клиринговые, по~диспутам).

\subsection*{Система управления картой (CMS)}

Система управления картой (Card Management System, CMS) ведёт жизненный цикл
карты: выпуск, смена статуса, блокировка, перевыпуск. CMS хранит реквизиты
и~связь со~счётом, поддерживает статусы и~лимиты, формирует задания
на~персонализацию и~фиксирует события, которые позже потребуются поддержке,
антифроду и~разбору диспутов.

На~этапе персонализации на~карту загружаются апплет EMV
(Europay/Mastercard/Visa, спецификация чиповых карт), ключи эмитента IMK
(Issuer Master Key, мастер-ключ эмитента), параметры управления риском эмитента
(IAC, Issuer Action Codes, и~прикладные офлайн-лимиты~--- последовательных
операций и~кумулятивной суммы), а~также данные платёжного приложения AID
(Application Identifier, идентификатор приложения) и~Application Label;
подробнее в~главе~\ref{ch:emv}~\cite{emvco-book3,emvco-book4}. Для~карт Мир
персонализация дополнительно включает установку приложения Мир (MPA, MIR
Payment Application) с~поддержкой криптографии ГОСТ; конкретный
криптографический профиль зависит от~программы эмиссии и~продолжающегося
перехода НСПК на~отечественные алгоритмы;
см.~раздел~\ref{sec:mir-certification}.

\subsection*{Авторизационный хост}

Авторизационный хост решает, можно ли по~карте провести конкретную операцию.
Когда запрос доходит от~карточной сети до~эмитента, хост последовательно
проверяет статус карты и~срок её действия, доступность средств или свободного
кредитного лимита, лимиты и~признаки нетипичного поведения,
см.~главу~\ref{ch:fraud}. Для~транзакций EMV он также верифицирует
криптограмму ARQC (Authorization Request Cryptogram, запросная криптограмма)
через HSM (Hardware Security Module, аппаратный
криптомодуль)~\cite{emvco-book2,emvco-book3}.

Если все проверки пройдены, хост формирует ответ \texttt{00} Approved
и~генерирует ARPC (Authorization Response Cryptogram), который возвращается
карте для взаимной аутентификации. Если хотя бы одна проверка не~проходит, хост
возвращает соответствующий response code, например \texttt{51} Insufficient
Funds, \texttt{05} Do Not Honor или \texttt{14} Invalid Card Number,
и~отклоняет транзакцию~\cite{emvco-book3}.

Отдельный сценарий~--- CMS недоступна, авторизационный хост работает. Без CMS
хост не~видит ни актуального статуса карты (заблокирована ли она, перевыпущена
ли), ни текущего лимита, и~у~него остаётся два варианта. Первый~--- отклонить
все транзакции; безопасно, но парализует приём по~всему портфелю. Второй~---
работать по~кэшу; хост хранит копию критических полей (статус, лимит, порог
офлайн-авторизации) со~сроком жизни записи кэша 5--15~мин и~продолжает
авторизовать в~пределах кэшированных лимитов. Транзакции сверх кэшированного
лимита хост отклоняет; по~остальным эмитент принимает риск одобрения
по~устаревшим данным. После восстановления CMS хост синхронизирует остатки
и~фиксирует расхождения. Для~портфеля в~миллионы карт разница между 5~мин
и~30~мин простоя CMS может означать сотни одобренных транзакций
по~заблокированным картам~\cite{emvco-book3}.

\subsection*{Сценарии эмитента (issuer scripting)}

Вместе с~ответом эмитент может прислать команду для карты~--- обновить
офлайн-PIN, сбросить счётчики, заблокировать приложение, изменить лимиты.
Механизм issuer scripting и~правила EMV для него разобраны
в~разделе~\ref{sec:emv-cryptogram}.

\subsection*{ACS~--- сервер аутентификации}

В~интернет-платежах одной проверки карты, лимита и~криптограммы недостаточно:
нужно ещё понять, действительно ли покупку инициировал держатель. У~эмитента
эту задачу решает ACS (Access Control Server)~--- компонент 3-D Secure
(протокол подтверждения держателя в~CNP-операциях), выбирающий между бесшовным
(\textit{frictionless}) сценарием и~дополнительным подтверждением по~логике RBA
(Risk-Based Authentication, аутентификация на~основе оценки риска);
см.~главу~\ref{ch:3ds}.

\needspace{6\baselineskip}
\subsection*{Управление токенами}

После выпуска карты её реквизиты живут на~пластике, в~телефоне, в~браузере
и~у~продавца в~виде сетевых токенов DPAN (Device PAN, токенизированный номер
карты для устройства); см.~раздел~\ref{sec:token-anatomy}. Эмитент участвует
в~этом жизненном цикле с~самого начала: получает от~провайдера токенизации TSP
(Token Service Provider, провайдер услуг токенизации) запрос на~токенизацию,
принимает решение об~одобрении после процедуры ID\&V (Identification \&
Verification, идентификация и~верификация) держателя, а~затем управляет
состоянием токена~--- активирует, приостанавливает или
удаляет~\cite{emvco-payment-tokenisation}.\footnote{\cite{visa-vts,nspk-tsp}.}

Токен остаётся связанным с~исходной картой. Если держатель блокирует карту,
эмитент обязан приостановить и~связанные с~ней токены. Если карта перевыпущена,
нужно обновить привязку токенов к~новому PAN (Primary Account Number, основной
номер счёта; см.~главу~\ref{ch:tokenization}), иначе держатель увидит
в~приложении банка новую карту, а~у~кошелька или продавца останется старая
привязка~\cite{emvco-payment-tokenisation,visa-vts}.

\subsection*{Работа с~диспутами}

Держатель карты с~претензией обращается не~в~карточную сеть и~не к~эквайеру,
а~к~эмитенту~--- своему банку: \enquote{я~не совершал эту покупку} или
\enquote{товар не~получен}. Эмитент проверяет обоснованность жалобы и~при
наличии оснований инициирует чарджбэк (chargeback, принудительный возврат
  средств держателю по~правилам карточной сети с~переносом ответственности
на~эквайера и~ТСП) через карточную сеть, направляя его эквайеру; жизненный цикл
диспута разобран
в~разделе~\ref{sec:disputes-lifecycle}~\cite{visa-core-rules}.\footnote{\cite{mc-chargeback-guide,visa-vcr-merchants}.}

Приняв обращение и~отнеся его к~подходящему классу, эмитент сопоставляет
ситуацию с~правилами сети: есть ли подходящий код причины, не~пропущены ли
сроки, достаточно ли документов и~соблюдены ли обязательные условия для
чарджбэка. После этого эмитент формирует чарджбэк-сообщение, а~если эквайер
присылает повторное представление с~\textit{compelling evidence}
(компенсирующие свидетельства~--- набор документов, подтверждающих легитимность
операции), эмитент решает, принять объяснение или идти дальше,
к~\textit{pre-arbitration} (досудебная стадия спора между эмитентом и~эквайером
перед арбитражем сети)~\cite{visa-core-rules}.\footnote{\cite{mc-chargeback-guide,visa-vcr-merchants}.}

\needspace{6\baselineskip}
\section{Эквайер}\label{sec:ia-acquirer}

\highlight{Эквайер (acquirer) подключает терминал или платёжный шлюз,
  доставляет авторизационный запрос в~карточную сеть, формирует и~отправляет
  клиринговый файл, принимает расчёт от~эмитентов и~перечисляет деньги продавцу
за~вычетом комиссий}. Если эмитент несёт кредитный риск держателя, то
эквайер~--- риск ТСП: расчётный, репутационный, правовой.

\practicenote{%
  Расчётный риск~--- риск того, что расчёт в~системе перевода средств
  не~состоится в~ожидаемое время; включает кредитную и~ликвидную
  компоненты~\cite{bis-cpmi-glossary}.

  Правовой риск~--- риск убытков из-за нарушения договоров и~допущенных правовых
  ошибок, несовершенства правовой системы, нарушения контрагентами нормативных
  актов или~различия юрисдикций у~филиалов и~контрагентов (п.~3.3 Указания Банка
    России от~15.04.2015~№~3624-У \enquote{О~системе управления рисками и~капиталом
  кредитной организации}~\cite{cbr-3624u}).

  Репутационный риск (риск потери деловой репутации)~--- риск убытков
  от~ухудшения восприятия банка клиентами, контрагентами, регулятором
  и~участниками рынка (Письмо Банка России от~30.06.2005~№~92-Т
    \enquote{Об~организации управления правовым риском и~риском потери деловой
репутации}~\cite{cbr-92t}). }

\subsection*{Система управления терминалами (TMS)}

В~офлайн-эквайринге первое физическое звено~--- терминал. У~крупного эквайера
TMS управляет десятками и~сотнями тысяч устройств. Главная задача~--- держать
на~каждом устройстве актуальную конфигурацию: обновлённые таблицы BIN (Bank
  Identification Number, банковский идентификационный номер;
см.~главу~\ref{ch:card-data}), свежие ключи, корректные параметры EMV.

Криптографические ключи терминала~--- ими он шифрует вводимый держателем PIN
и~подписывает сообщения, отправляемые эквайеру; устройство и~иерархия этих
ключей разобраны в~разделе~\ref{sec:crypto-key-hierarchy}. Исторически загрузка
ключей на~терминал шла очной процедурой с~двумя участниками в~защищённой зоне
эквайера; на~парке в~тысячи устройств она уже не~справлялась. Удалённая
инъекция ключей (Remote Key Injection, RKI) снимает это ограничение: терминалы
разворачиваются и~обновляются без физического
доступа~\cite{visa-tpa-registration-program-faqs}. Церемония при этом
сдвигается на~уровень корневых ключей RKI и~мастер-ключей HSM
(см.~раздел~\ref{sec:crypto-key-ceremony}); из~них через защищённый канал
производятся ключи отдельных терминалов.

\subsection*{Конфигурация ТСП и~параметры маршрутизации}

Для интернет-эквайринга и~омниканальных сценариев одних терминалов мало. Нужна
ещё надстройка, в~которой живут профили ТСП, конфигурация MID/TID (Merchant ID
  / Terminal ID, идентификаторы ТСП и~терминала в~учётных системах эквайера
и~карточной сети), правила дескрипторов, маршрутизация по~процессорам,
параметры 3-D Secure, лимиты на~возвраты и~выплаты, а~также связка с~платёжным
шлюзом из~главы~\ref{ch:gateway}.

Эта надстройка определяет, какой MCC и~какой дескриптор увидит эмитент, через
какой BIN эквайера пойдёт транзакция и~в~какую программу мониторинга попадёт
ТСП. Если такие решения живут в~разрозненных таблицах, тикетах и~ручных
настройках в~кабинете процессора, эквайер теряет управляемость: одна и~та же
конфигурация ТСП по-разному выглядит для сети, команды риска
и~финансовой сверки~\cite{visa-merchant-data-standards-manual}.

У~крупных эквайеров в~этой надстройке сходятся подключение ТСП, маршрутизация
между несколькими процессорами и~операционная политика: кому дать прямой
маршрут на~конкретного процессора, кому оставить один канал, а~кому включить
усиленный фиксированный резерв (\textit{holdback}) или ручную проверку выплат.

\subsection*{Подключение ТСП}

Прежде чем дать ТСП доступ к~карточной сети, эквайер должен установить, кто
перед ним, какой бизнес он ведёт и~какой риск несёт. Подключение ТСП~--- набор
проверок и~настроек, которые задают допуск и~экономику отношений.

При идентификации ТСП (KYB, Know Your Business) проверяют юридическое лицо,
бенефициаров и~санкционные списки. При классификации MCC ТСП присваивают код
категории, от~которого зависят ставки interchange, правила 3DS и~допустимость
операции. При \textit{underwriting} оценивают финансовый риск и~задают ставку
комиссии, размер фиксированного резерва (\textit{holdback}) и~плавающего
резерва (\textit{rolling reserve}). Разбор каждого шага~---
в~разделе~\ref{sec:kyb-merchants}.

\subsection*{Мониторинг ТСП}

Эквайер непрерывно наблюдает за~долей чарджбэков, долей фрода, долей возвратов
и~аномалиями частоты (\textit{velocity}~--- всплески числа операций по~карте
или ТСП относительно ожидаемого профиля). Превышение порогов карточных схем
ведёт к~штрафам и~усиленному контролю; пропущенный \textit{bust-out} (схема,
  при которой ТСП за~короткий срок собирает большой объём авторизаций
  и~выводит средства
до~истечения срока, в~течение которого допустим чарджбэк) обходится эквайеру
в~миллионы. \highlight{Ответственность за~ТСП несёт эквайер}: если ТСП оказался
мошенником или нарушил правила карточной сети, штрафы и~финансовые потери
остаются у~эквайера. Метрики мониторинга и~обмен с~ФинЦЕРТ (Центр мониторинга
  и~реагирования на~компьютерные атаки в~кредитно-финансовой сфере Банка
России)~\cite{cbr-fincert-fraud-overview} разобраны
в~разделе~\ref{sec:merchant-monitoring}, мониторинговые программы сетей~---
VAMP (Visa Acquirer Monitoring Program) и~MC~ECP (Mastercard Excessive
Chargeback Program)~---
в~разделе~\ref{sec:disputes-monitoring}~\cite{visa-vamp-intro,mc-tpr}.

\needspace{6\baselineskip}
\subsection*{Ответственность за расчёт}

Для продавца эквайринг заканчивается в~момент поступления денег. Эквайер
собирает клиринг, удерживает комиссии, применяет плавающий резерв
(\textit{rolling reserve}) и~перечисляет средства продавцу.

Эквайер обязан показать состав каждой выплаты: какая сумма удержана как
комиссия, какая отправлена в~резерв, какая зачтена против возвратов и~в~какой
валюте считается обязательство. Кроме того, на~нём остаётся риск \textit{late
presentment} (просрочки представления операции в~клиринг сверх установленного
сетью окна; подробнее см.~раздел~\ref{sec:tx-clearing}), временных разрывов
по~чарджбэкам и~кассовых разрывов между получением денег от~сети и~выплатой
продавцу. Эквайер же задаёт политику выплат в~нотации~T+N: T~--- день
совершения операции (\textit{trade date}), N~--- сдвиг в~рабочих банковских
днях до~зачисления на~счёт ТСП. T+0~--- в~тот~же день, T+1~--- на~следующий
рабочий день, T+2~--- через два; еженедельные выплаты и~отложенный расчёт
для~высокорисковых MCC~--- варианты той~же шкалы. Сутки от~T отсекаются
по~\textit{cut-off} эквайера: операция после отсечки относится к~следующему~T.

На~счёт ТСП приходит нетто после удержаний на~стадиях жизненного цикла операции
(авторизация, клиринг, межбанковский расчёт~--- разбор
в~гл.~\ref{ch:tx-lifecycle}): торговая уступка (\textit{merchant discount
  rate}, MDR, синоним \textit{merchant service charge}, MSC; разложение
в~разделе~\ref{sec:ecosystem-interchange}), отчисления в~плавающий
и~фиксированный резервы (разбор~--- в~следующем абзаце), зачёты против
возвратов и~чарджбэков. Расчёт между сетью и~банком-эквайером
(см.~раздел~\ref{sec:tx-settlement}) идёт по~своему графику. Когда эквайер
выплачивает ТСП быстрее, чем сеть рассчитывается с~ним (мгновенные выплаты, T+0
на~рынках, где сетевой расчёт T+1), разницу эквайер покрывает собственными
средствами~--- короткий кредит ТСП на~несколько дней. Стоимость такого
кредитования заложена в~наценку эквайера в~составе MSC и~повышает тариф
для~быстрых выплат.

Фиксированный и~плавающий резервы нужны из-за временного разрыва между
авторизацией и~истечением срока, в~течение которого держатель ещё вправе
оспорить операцию; от~авторизации до~этого момента может пройти несколько
месяцев. Если продавец за~это время прекратил деятельность или не~может покрыть
возвраты, потери падают на~эквайера. Резерв~--- буфер для таких потерь.
Его размер отражает оценку риска ТСП: для продавца с~высокой долей чарджбэков,
сезонным бизнесом или без истории эквайер устанавливает более высокий процент
удержания.\footnote{Эти резервы эквайера~--- holdback и~rolling reserve~---
  удерживаются на~стороне эквайера. Не~путать с~резервами в~продуктовом реестре
  платформы (refund reserve, minimum balance, holdback~---
  см.~раздел~\ref{sec:ledger-balances}): там та~же терминология описывает
  удержание у~продавца на~стороне маркетплейса\slash PSP, отдельным механизмом
и~по~собственным правилам.}

\subsection*{Завершение работы с~ТСП}

После расторжения договора риск эквайера сохраняется на~те~же месяцы, ради
которых создан резерв. ТСП расторгает договор сам или эквайер отключает его
из-за превышения порогов мониторинга, нарушения правил схемы или подозрения
на~\textit{bust-out}; в~обоих случаях операции последних месяцев ещё можно
оспаривать, и~за~них отвечает эквайер.

Резерв поэтому возвращают не~в~день расторжения, а~по~мере того, как
истекает срок оспаривания накопленных операций; высвобождение растягивается
до~полугода после последней транзакции (привязка резерва к~сроку
оспаривания разобрана в~гл.~\ref{ch:disputes}). Перед возвратом эквайер
зачитывает против резерва незакрытые чарджбэки, штрафы схемы и~неоплаченные
комиссии~--- продавец получает остаток.

Ответственность сохраняется и~в~процедурах схемы: представление возражений
(\textit{representment}) и~арбитраж по~операциям, совершённым до~расторжения,
проходят по~срокам схемы уже после отключения ТСП. Эквайер откладывает
финальный расчёт с~продавцом до~закрытия этих диспутов. Само отключение
затрагивает и~учётные системы эквайера: MID и~TID закрывают в~CMS, терминалы
снимают с~обслуживания в~TMS, маршруты и~дескрипторы исключают из~конфигурации.

Если договор расторгнут из-за риска, эквайер вносит продавца в~реестр
расторгнутых ТСП своей схемы~--- исторически такой реестр называли Terminated
Merchant File (TMF). У~Mastercard это теперь MATCH (Mastercard Alert to
Control High-risk Merchants): запись содержит код причины~--- компрометация
данных, отмывание, избыточные чарджбэки, избыточный фрод, осуждение
за~мошенничество~--- и~хранится пять лет~\cite{mc-spme-2025}. Правила Visa
предписывают эквайеру участие в~Visa Merchant Screening Service
(VMSS)~\cite{visa-merchant-screening-service}. Реестры схем раздельны: MATCH
стал отраслевым стандартом де-факто, и~эквайеры сверяются с~ним и~по~картам
Visa, хотя обязывает к~нему только Mastercard. Следующий эквайер видит запись
при~андеррайтинге (см.~раздел~\ref{sec:kyb-merchants}) и~учитывает расторжение
у~предыдущего эквайера как риск.

В~России отдельного публичного реестра расторгнутых ТСП у~платёжной системы
Мир нет: программа безопасности НСПК требует от~эквайера подтверждать
соответствие ТСП требованиям защиты данных и~предусматривает штрафы
за~нарушения, но~общего списка недобросовестных продавцов НСПК
не~ведёт~\cite{nspk-mir-security}. Допуск и~отключение ТСП опираются здесь
на~договор эквайера с~ТСП и~на~инструменты Банка России~--- платформу
\enquote{Знай своего клиента} (ЗСК), которая присваивает юрлицам и~ИП уровень
риска по~115-ФЗ (см.~раздел~\ref{sec:kyc-zsk}).

\section{Платёжный фасилитатор (PayFac) и~модель субмерчантов}\label{sec:ia-payfac}

Классическая модель эквайринга работает, пока ТСП немного и~каждое подключено
как отдельный клиент банка. Для маркетплейса или SaaS-платформы с~тысячами
контрагентов банковское подключение становится узким местом: оформление каждого
продавца занимает дни или недели, а~бизнес-модель предполагает подключение
за~минуты.

\highlight{В~модели PayFac платформа регистрируется спонсирующим эквайером
  в~карточной сети как PayFac и~получает право подключать sponsored merchants
по~ускоренной схеме}. У~субмерчанта может быть собственный идентификатор, но
договор и~расчёты идут через PayFac и~его эквайера. Эквайер перечисляет
средства PayFac, а~тот уже рассчитывается с~субмерчантами. Нового продавца
подключают за~минуты вместо недель: PayFac автоматизирует идентификацию
субмерчанта, присваивает MCC и~начинает приём платежей, не~дожидаясь полного
банковского
цикла~\cite{visa-acceptance-entities,visa-payment-facilitator-marketplace-risk-guide}.

Ответственность за~субмерчантов ложится на~PayFac. Перед карточной
сетью и~эквайером он отвечает за~их действия как за~свои: мошенничество
субмерчанта, рост чарджбэков, нарушение правил сети~--- всё это штрафы
и~потери PayFac. Поэтому он строит собственные процессы идентификации,
мониторинга ТСП и~работы с~диспутами, по~строгости мало отличающиеся от
банковских~\cite{visa-acceptance-entities,visa-payment-facilitator-marketplace-risk-guide}.

Карточная сеть видит PayFac как единый субъект; что происходит внутри его
портфеля субмерчантов, ей напрямую не~известно. Нарушение правил субмерчантом
сеть замечает только через отклонения в~агрегированных показателях PayFac.

\importantnote{%
  PayFac и~ISO (Independent Sales Organization)~--- разные роли.
  ISO привлекает и~обслуживает ТСП для банка-эквайера, получая
  агентское вознаграждение без прямого участия в~расчёте. PayFac
  находится в~денежном потоке: эквайер перечисляет средства ему,
  а~он затем рассчитывается с~субмерчантами~\cite{visa-tpa-registration-program-faqs,visa-acceptance-entities}.%
}

\subsection*{Когда выбирать PayFac}

Платформа, выводящая на~приём множество мелких продавцов, выбирает одну из~трёх
моделей. \textbf{PayFac} (Payment Facilitator)~--- платформа регистрируется
в~сети через спонсирующего эквайера, ведёт идентификацию субмерчантов, несёт
ответственность за~их чарджбэки и~контролирует выплаты; оправдан при сотнях
субмерчантов и~потребности в~быстром подключении. \textbf{MoR} (Merchant of
Record)~--- платформа выступает единственным ТСП в~глазах сети, а~субмерчанты
остаются её внутренней абстракцией; проще регуляторно, но налоговые
обязательства, чарджбэки и~возвраты замыкаются на~платформу; подходит для SaaS
с~единым продуктом. \textbf{pass-through}, прямой эквайринг,~--- каждое ТСП
заключает договор с~эквайером напрямую, а~платформа только маршрутизирует
трафик; ответственность платформы минимальна, выплатами она не~управляет,
подключение каждого ТСП занимает недели~\cite{visa-acceptance-entities}.

Быстрое подключение в~модели PayFac даёт платформе обязанности перед сетью
и~эквайером:
\begin{itemize}
  \item ответственность за~мошенничество и~чарджбэки субмерчантов перед
    карточной сетью и~спонсирующим эквайером;
  \item регистрация через эквайера как Third-Party Agent в~Visa (Program
    Request Management) и~аналогичные процедуры
    Mastercard~\cite{visa-payment-facilitator-marketplace-risk-guide};
  \item для платежей в~CEMEA~--- сертификация Visa Ready Acceptance
    и~подписание Payment Facilitator Agreement
    с~Visa~\cite{visa-payment-facilitator-marketplace-risk-guide};
  \item собственный процесс идентификации субмерчантов и~ПОД/ФТ
    (противодействие отмыванию доходов и~финансированию терроризма)~---
    эквайер требует от~PayFac выполнять KYB и~санкционную проверку по~Visa
    Global Acquirer Risk Standards.
\end{itemize}

\practicenote{%
  Если вы строите платформу и~выбираете между PayFac и~прямым
  эквайрингом для каждого ТСП, оцените, сколько ТСП реально
  планируете подключить в~первый год. На~десятках прямой эквайринг
  обычно проще и~дешевле. На~сотнях и~тысячах PayFac часто окупает себя
  скоростью подключения, но требует серьёзных инвестиций
  в~риск, соблюдение требований и~повседневное обслуживание.%
}

\subsection*{Параллель в~России~--- банковский платёжный агент}

PayFac-модель родилась внутри правил карточных сетей и~применяется там же.
Местный правопорядок оперирует другой конструкцией~--- институтом банковского
платёжного агента (БПА) и~банковского платёжного субагента, закреплённым
статьями~14 и~14.2 161-ФЗ~\cite{fz-161}; для БПА, осуществляющих операции
платёжного агрегатора, дополнительно применяется Указание Банка России
от~06.04.2020~№~5429-У \enquote{О~ведении перечня
БПА-агрегаторов}~\cite{cbr-5429u}. Кредитная организация вправе привлекать БПА
для приёма наличных, переводов, идентификации и~иных операций в~случаях
и~в~пределах, установленных законом.

Для платформы, подключающей субмерчантов и~работающей с~местным эквайрингом,
БПА~--- отдельная правовая конструкция со~своими обязанностями; прямой заменой
PayFac из~правил Visa/Mastercard она не~служит. Для БПА-операторов платёжного
агрегатора предусмотрено включение в~перечень, который ведёт Банк России;
общего публичного реестра всех БПА не~существует. К~обязанностям БПА относятся
выполнение требований 115-ФЗ по~идентификации и~контролю операций
(см.~главы~\ref{ch:kyc-kyb} и~\ref{ch:aml}), ведение специального банковского
счёта, на~который зачисляются принимаемые средства и~с~которого они
перечисляются, и~соблюдение ограничений по~перечню операций, которые БПА имеет
право выполнять от~имени банка. Банковский платёжный субагент находится
на~следующем уровне договорных отношений и~привлекается уже самим БПА
с~дополнительными ограничениями.

У~платформы выбор идёт не~между \textit{Stripe Connect} и~\textit{Adyen for
Platforms}, а~между тремя вариантами на~местном рынке: договор БПА с~одним или
несколькими банками; интеграция через PSP, который сам уже имеет статус БПА
и~предлагает решение для маркетплейсов как собственный продукт; роль
спонсируемого PayFac у~эквайера~--- участника платёжной системы \enquote{Мир}.
Архитектуру выплат и~сверки проектируют под выбранную модель; ссылка
на~универсальные правила PayFac без оговорки про БПА неточна.

Прямой договор БПА с~банком даёт контроль над приёмом и~клиентским путём,
но~обязывает платформу самостоятельно вести специальный банковский счёт,
поддерживать включение в~перечень БПА-агрегаторов Банка России и~выполнять
требования 115-ФЗ по~идентификации субмерчантов (юридических лиц и~ИП).
Идентификация клиентов-физлиц к~функциям БПА не~относится: Федеральный закон
от~20.02.2026~№~44-ФЗ \enquote{Об~исключении БПА из~идентификации физлиц
по~115-ФЗ}~\cite{fz-44} изъял у~банковских платёжных агентов право проводить
идентификацию (в~том числе упрощённую) и~обновление сведений о~физических
лицах\footnote{Вступил в~силу 03.03.2026; существующие договоры подлежат
приведению в~соответствие в~течение 180~дней.}. Интеграция через PSP-БПА (PSP
в~статусе банковского платёжного агента) снимает с~платформы прямые
обязательства перед регулятором: специальный счёт, перечень и~процедуры
идентификации остаются на~PSP. Цена этого варианта~--- маржа и~меньшая
гибкость: параметры приёма, тарифы и~сценарии возврата определены продуктом
PSP, и~собственный сценарий маркетплейса чаще приходится укладывать в~его API.
Спонсируемый PayFac у~эквайера \enquote{Мира} ближе всего к~привычной
международной модели и~даёт быструю регистрацию субмерчантов внутри карточной
сети, но~работает только для карточных операций; СБП-приём и~переводы по~номеру
счёта остаются за~пределами этого договора и~подключаются отдельно. Сверка
и~казначейство платформы поэтому делятся на~карточный и~некарточный потоки
с~самого старта проекта.
